\chapter{Reconstructing CMC surfaces from harmonic maps}\label{sec:cmc-reconstruction}
The construction in \cite{Hel01}, Chapter 2, pp.~19--21, turns a harmonic sphere-valued map into CMC surfaces. We derive the formulas with the convention $h(X,Y)=\langle D_X\nu,Y\rangle$ and trace mean curvature. The period condition and the possible failure of immersion are part of the statement.

\begin{lemma}[The closed reconstruction form]\label{lem:cmc-reconstruction-closed-form}
Let $\Omega\subset\mathbb R^2$ be open and let $u:\Omega\to S^2$ be smooth and harmonic. The $\mathbb R^3$-valued one-form
\[
 \alpha=(u\times u_y)\,dx-(u\times u_x)\,dy
\]
is closed. On a simply connected domain there is a smooth map $B:\Omega\to\mathbb R^3$ with $dB=\alpha$, unique up to translation. It satisfies $\Delta B=2u_x\times u_y$.
\end{lemma}
\begin{proof}
The coefficient of $dx\wedge dy$ in $d\alpha$ is $-u\times(u_{xx}+u_{yy})$, which vanishes because $\Delta u=-|Du|^2u$. On a simply connected domain a closed form has a primitive: integrate along a path from a fixed base point; Stokes' theorem applied to a homotopy shows path independence, component by component. Differentiating this integral gives the form, and any two primitives have zero derivative and differ by a constant. Finally
\[
 B_{xx}+B_{yy}=u_x\times u_y-u_y\times u_x+u\times(u_{yx}-u_{xy})=2u_x\times u_y.
\]
\end{proof}

\begin{proposition}[The global period obstruction]\label{prop:cmc-reconstruction-periods}
On a connected, not necessarily simply connected, planar domain, the form in Lemma \ref{lem:cmc-reconstruction-closed-form} has a single-valued primitive if and only if its integral around every closed piecewise smooth curve is zero. Without that condition the construction still works locally and on the universal cover.
\end{proposition}
\begin{proof}
The integral of $dB$ around a closed curve is zero by the fundamental theorem of calculus. Conversely, vanishing periods make the path integral from a base point independent of the chosen path: traverse one path and then the reverse of the other. Local differentiation produces a smooth primitive. On the universal cover, the pulled-back form is closed and the domain is simply connected, so the preceding lemma applies in local coordinates. Different lifts of a point can give values of $B$ differing by a period vector. Thus local solvability alone does not imply a single-valued map on the original domain.
\end{proof}

\begin{lemma}[The algebra of the two reconstructed maps]\label{lem:cmc-reconstruction-algebra}
Assume $dB=\alpha$ and set $X_\varepsilon=B+\varepsilon u$ for $\varepsilon\in\{1,-1\}$. Write
\[
 E=|u_x|^2+|u_y|^2,\quad d=u\cdot(u_x\times u_y),\quad
 \lambda_\varepsilon^2=E-2\varepsilon d.
\]
Then $\lambda_\varepsilon^2\ge0$, and
\[
 |(X_\varepsilon)_x|^2=|(X_\varepsilon)_y|^2=\lambda_\varepsilon^2,\qquad
 (X_\varepsilon)_x\cdot(X_\varepsilon)_y=0,
\]
\[
 (X_\varepsilon)_x\times(X_\varepsilon)_y=-\varepsilon\lambda_\varepsilon^2u,
 \qquad \Delta X_\varepsilon=-\varepsilon\lambda_\varepsilon^2u.
\]
\end{lemma}
\begin{proof}
Put $a=u_x$, $b=u_y$ and $Jv=u\times v$ for vectors tangent to $S^2$ at $u$. Then $J^2=-I$, $|Jv|=|v|$, and
\[
 (X_\varepsilon)_x=Jb+\varepsilon a,\qquad
 (X_\varepsilon)_y=-Ja+\varepsilon b=-\varepsilon J(X_\varepsilon)_x.
\]
Since $a\cdot Jb=-d$, the first squared norm is $E-2\varepsilon d$; the second and the orthogonality follow from the last relation. For a tangent vector $v$, $v\times Jv=|v|^2u$, proving the cross-product formula. The previous lemma and harmonicity give $\Delta X_\varepsilon=2a\times b-\varepsilon Eu=(2d-\varepsilon E)u$, as claimed.
\end{proof}

\begin{theorem}[CMC reconstruction on the immersion locus]\label{thm:cmc-harmonic-reconstruction}
Under the preceding hypotheses, $X_\varepsilon$ is a conformal immersion on the open set $\{\lambda_\varepsilon^2>0\}$. Its induced metric is $\lambda_\varepsilon^2(dx^2+dy^2)$, and its trace scalar mean curvature relative to the chosen normal $u$ is $H=\varepsilon$. Its mean curvature vector is $-\varepsilon u$. At a point where $\lambda_\varepsilon=0$, its entire differential vanishes; no immersion assertion is made there. This is the reconstruction of \cite{Hel01}, pp.~19--21, with signs and normalization made explicit.
\end{theorem}
\begin{proof}
The positive conformal factor gives a rank-two differential by Lemma \ref{lem:cmc-reconstruction-algebra}. Since both derivatives are perpendicular to $u$, this is a unit normal. For the induced metric, $\Delta_g X_\varepsilon=\lambda_\varepsilon^{-2}\Delta X_\varepsilon=-\varepsilon u$. The immersion identity $\Delta_gX=\vec H=-Hu$ yields $H=\varepsilon$. The parameter orientation instead gives unit normal $-\varepsilon u$, relative to which the scalar trace is $-1$. At a zero of the conformal factor, both coordinate derivatives have norm zero. Thus the formula can produce a degenerate map, even a constant map, as well as an immersion.
\end{proof}

\begin{proposition}[Hopf differential and degeneracy of reconstruction]\label{prop:cmc-reconstruction-hopf}
Let $q=\langle u_z,u_z\rangle$ be the coefficient of the Hopf differential. Then
\[
 \lambda_+^2\lambda_-^2=16|q|^2.
\]
If $\Omega$ is connected and $q$ is not identically zero, both reconstructions are immersions away from a discrete set. At a zero of $q$, at least one reconstruction has zero differential. A conformal harmonic map can therefore make one reconstruction degenerate everywhere.
\end{proposition}
\begin{proof}
Since $a,b$ lie in the plane perpendicular to $u$, $d^2=|a|^2|b|^2-(a\cdot b)^2$. Consequently
\[
 (E-2d)(E+2d)=(|a|^2-|b|^2)^2+4(a\cdot b)^2=16|q|^2.
\]
Both factors are nonnegative. The Hopf coefficient is holomorphic by Proposition \ref{prop:harmonic-hopf-differential}, so its zeros are discrete unless it vanishes identically, by the identity theorem for holomorphic functions. At a zero the product vanishes, and the differential of the corresponding reconstruction is zero. The last possibility is realized explicitly in Example \ref{ex:stereographic-cmc-reconstruction}.
\end{proof}
